% from ega0-16-fr.tex
\oldpage[0\textsubscript{IV-1}]{39}

\begin{corollary}[16.5.11]
\label{0.16.5.11-fr}
Soit $A$ un anneau local noethérien~; supposons qu'il existe un $A$-module de
type fini $M$ de support $\operatorname{Spec}(A)$, qui soit un module de
Cohen-Macaulay. Alors, pour tout idéal premier $\mathfrak{p}$ de $A$ on a
\begin{equation}
\label{0.16.5.11.1-fr}
\tag{16.5.11.1}
  \dim(A)=\dim(A/\mathfrak{p})+\dim(A_{\mathfrak{p}}).
\end{equation}
\end{corollary}

\begin{proof}
En effet, $\operatorname{Supp}(M/\mathfrak{p}M)
=\operatorname{Spec}(A/\mathfrak{p})$ ($0_{\mathrm I}$, 1.7.5) et
$\operatorname{Supp}(M_{\mathfrak{p}})
=\operatorname{Spec}(A_{\mathfrak{p}})$~; la relation (16.5.11.1) est donc un
cas particulier de (16.5.10.1).
\end{proof}

\begin{corollary}[16.5.12]
\label{0.16.5.12-fr}
Tout anneau quotient d'un anneau local noethérien $A$ vérifiant les conditions
de (16.5.11) (en particulier tout anneau quotient d'un anneau local de
Cohen-Macaulay) est \emph{caténaire}.
\end{corollary}

\begin{proof}
Il suffit évidemment de prouver que $A$ lui-même est caténaire c'est-à-dire
que, pour deux idéaux premiers $\mathfrak{p}\subset\mathfrak{q}$ de $A$, on a
(16.1.4.2)
\[
  \dim(A_{\mathfrak{q}})
  =\dim(A_{\mathfrak{p}})
  +\dim(A_{\mathfrak{q}}/\mathfrak{p}A_{\mathfrak{q}}).
\]
Or, en vertu de (16.5.10, (i)) $A_{\mathfrak{q}}$ vérifie les mêmes
hypothèses que $A$, et la relation précédente n'est donc autre que
(16.5.11.1) appliquée à $A_{\mathfrak{q}}$ et à l'idéal premier
$\mathfrak{p}A_{\mathfrak{q}}$ de $A_{\mathfrak{q}}$.
\end{proof}

\begin{env}[16.5.13]
\label{0.16.5.13-fr}
Soient $A$ un anneau noethérien, $M$ un $A$-module de type fini. On dit que
$M$ est un \emph{$A$-module de Cohen-Macaulay} si, pour tout idéal premier
$\mathfrak{p}$ de $A$, $M_{\mathfrak{p}}$ est un $A_{\mathfrak{p}}$-module
de Cohen-Macaulay~; en vertu de (16.5.10) cette définition coïncide avec
(16.5.1) lorsque $A$ est \emph{local}. On dit que $A$ est un
\emph{anneau de Cohen-Macaulay} si tous les $A_{\mathfrak{p}}$ le sont. Il est
clair que si $M$ (resp. $A$) est un $A$-module de Cohen-Macaulay (resp. un
anneau de Cohen-Macaulay), $S^{-1}M$ est un $S^{-1}A$-module de Cohen-Macaulay
(resp. $S^{-1}A$ est un anneau de Cohen-Macaulay) pour toute partie
multiplicative $S$ de $A$.
\end{env}

% (continued in ega0-17-fr.tex, before its first page marker)
\section{Anneaux réguliers}
\label{section:0.17-fr}

\subsection{Définition des anneaux réguliers}
\label{subsection:0.17.1-fr}

\begin{proposition}[17.1.1]
\label{0.17.1.1-fr}
Soient $A$ un anneau local noethérien de dimension $n$, $\mathfrak{m}$ son
idéal maximal, $k=A/\mathfrak{m}$ son corps résiduel. Les conditions suivantes
sont équivalentes~:
\begin{enumerate}
  \item[\rm a)] L'homomorphisme canonique surjectif de $k$-modules gradués
  (15.1.1.1)
  \[
    \varphi:\mathbf{S}^*_k(\mathfrak{m}/\mathfrak{m}^2)
    \longrightarrow\operatorname{gr}^*_{\mathfrak{m}}(A)
  \]
  (où le second membre est le $k$-module gradué associé à $A$ muni de la
  filtration $\mathfrak{m}$-préadique) est bijectif.
  \item[\rm b)] On a
  $\operatorname{rg}_k(\mathfrak{m}/\mathfrak{m}^2)=n$.
  \item[\rm c)] L'idéal $\mathfrak{m}$ admet un système de générateurs de $n$
  éléments.
  \item[\rm d)] L'idéal $\mathfrak{m}$ admet un système de générateurs qui est
  une suite $A$-régulière.
\end{enumerate}
\end{proposition}


\begin{proof}
En vertu du lemme de Nakayama,
$\operatorname{rg}_k(\mathfrak{m}/\mathfrak{m}^2)$ est le plus petit nombre
d'éléments d'un système de générateurs de $\mathfrak{m}$, donc b) et c) sont
équivalentes. D'autre part, si $(x_i)_{1\leq i\leq r}$ est un système de
générateurs de $\mathfrak{m}$ dont les classes $\overline{x}_i$ mod.
$\mathfrak{m}^2$ forment une base du $k$-espace vectoriel
$\mathfrak{m}/\mathfrak{m}^2$, $\mathbf{S}^*(\mathfrak{m}/\mathfrak{m}^2)$
est isomorphe à l'anneau de
